\begin{table*}[t]
\centering
\caption{Screen stability when one input is removed.}
\label{tab:anomaly-feature-ablation}
\scriptsize
\setlength{\tabcolsep}{1pt}
\begin{tabularx}{\textwidth}{@{}l>{\raggedleft\arraybackslash}X>{\raggedleft\arraybackslash}X>{\raggedleft\arraybackslash}X>{\raggedleft\arraybackslash}X>{\raggedleft\arraybackslash}X@{}}
\toprule
Detector feature omitted & Flags & Case flags & Control flags & Enrichment & Overlap with main screen \\
\midrule
Self-citation rate & 53 & 49 (0.70\%) & 4 (0.04\%) & 16.14 & 0.963 \\
Coauthor-citation rate & 50 & 46 (0.66\%) & 4 (0.04\%) & 15.15 & 0.907 \\
Reciprocity & 52 & 48 (0.69\%) & 4 (0.04\%) & 15.81 & 0.981 \\
Local clustering & 55 & 51 (0.73\%) & 4 (0.04\%) & 16.80 & 0.964 \\
Outgoing HHI & 51 & 47 (0.68\%) & 4 (0.04\%) & 15.48 & 0.962 \\
\bottomrule
\end{tabularx}
\par\vspace{2pt}\begin{minipage}{0.96\linewidth}\scriptsize Each row removes one screen input. The cutoff and random seed stay fixed, and the Case and Control percentages use the same eligible rows. Overlap is the Jaccard share of flagged rows retained by both screens.
\end{minipage}
\end{table*}
